% !TeX root = ../papers/nc-p3.tex
\section{小標数の全有理通常倍率における正則性の持上げ}
\label{sec:rank-fifteen-small-residues}
\begin{theorem}\label{thm:rank-fifteen-small-residues}
$p=7,11,13,17,19$ のいずれかとし、標数零の幾何的な倍率 $(z,\lambda)$ がある付値で
$\mathbb F_p$ の通常倍率に還元されるとする。正の固定曲線のこの倍率の上にある
反対階数15の全窓は標準Koszul左右AS正則である。
各始点の経路の直和は四つの四次中心元の多項式環上で階数256の自由加群になる。
\end{theorem}
\begin{proof}
定理\ref{thm:rank-fifteen-root-exchange}と反対代数の対応により、
一つの倍率のある一根・一表示で四中心商の証明書があれば、両表示の全根が同じ証明書を持つ。
添字保持同型は四期の中心と商を運び、語の反転も四始点を並べ替えて商の次元を保つ。
倍率反転による二つの変換も使うと、剰余倍率では
\[
 (z,\lambda),\quad(z,1/(z\lambda)),\quad(1/z,1/\lambda),\quad(1/z,z\lambda)
\]
の組から一代表を選べばよい。後の二点は $z^2\ne1$ により前の二点と異なり、
$z\lambda^2\ne1$ により各対の中でも異なる。

\path{computations/verify_rank_fifteen_small_residues.py} は
$z\lambda(z^2-1)(\lambda^2-1)(z\lambda^2-1)(z^2\lambda^2-1)\ne0$ の全剰余対を
列挙し、この四点軌道に分ける。代表ごとに $P(n)$ の最小次数のモニック既約因子
$g$ を選び、剰余体 $\mathbb F_p[n]/(g)$ の根で第一表示を作る。
表の後三列は $\deg g$ による代表数である。
\[
\begin{array}{c|r|rrr}
 p&\text{通常倍率数}&1&2&4\\\hline
 7&4&0&0&1\\
 11&40&3&4&3\\
 13&72&6&6&6\\
 17&156&14&15&10\\
 19&208&12&20&20
\end{array}
\]
全120代表で始点 $0,1$、文字 $0,2$ の六次挿入行列は順に階数 $36,34$ を持つ。
続いて長さ4,5の経路次元 $35,56$、文字別の四次中心次元 $(3,2,0,0)$ を確認し、
中心から四式を選ぶ。全四始点でその商の長さ0から13の次元は
$1,4,10,20,31,40,44,40,31,20,10,4,1,0$ となる。

任意に指定した標数零の持上げで、モニックな $P$ の全根は付値環の整閉包に属し、
定数項は単元である。代数閉体まで延長して一次因子の積を還元すると、
選んだ剰余根に落ちる根が少なくとも一つある。重根でもこの議論は使える。
通常域のGram分母は各検査枠で可逆で、経路と中心の次元は一般点からの下界に等しい。
命題\ref{prop:fixed-residual-regularity}の持上げの証明により、
その標数零の根に全次数の自由性と左右Koszul完全性が従う。
最後に標数零で根交換・反転を施せばよい。このため剰余の四次式の分解型や
剰余での階数の低下を制限する必要はない。
\end{proof}
